\section{Benchmark Metrics}\label{sec:bench}
All measurements target the size-06 digest and the keyed-MD5 construction on a
host with 10 logical cores. Rates are in millions of keys hashed per second
(Mkeys/s).

\begin{table}[H]
\centering
\footnotesize
\begin{tabularx}{\linewidth}{X r r}
\hline
\textbf{Measurement} & \textbf{Aggregate} & \textbf{Per node/core}\\
\hline
BM1 single-core baseline            & 2.55 & 2.55\\
BM2 scaling, 8 processes            & 10.67 & 1.33\\
BM2 scaling, 10 processes           & 10.95 & 1.09\\
BM3 distributed, 8 TCP workers      & 12.56 & $\approx$1.75\\
\hline
\end{tabularx}
\caption{Hashing throughput (Mkeys/s). Full sweep in \texttt{results/benchmarks.txt}.}
\label{tab:bench}
\end{table}

\paragraph{CPU-bound vs.\ network-bound.} Instrumenting each worker to split its
wall time into hashing versus time blocked on the socket shows the workload is
overwhelmingly compute-bound: across the distributed run, aggregate compute time
was $69.1$\,s against $0.03$\,s of socket wait, i.e.\ \textbf{99.96\%} of time is
spent hashing. The coordination protocol is therefore not the bottleneck; total
throughput scales with available CPU cores until the host saturates
(\autoref{tab:bench}, where per-core rate falls once the process count exceeds
the physical core budget).

\paragraph{End-to-end.} The full recovery of \texttt{ykusrq} took $738$\,s
single-threaded versus $142$\,s on the contended cluster; an uncontended
$\sim$12.6\,Mkeys/s cluster covers the whole $2.18\times10^9$ space in well under
that.

% Optional figure: render results/benchmarks.json with notebooks/graphing.ipynb.
% \begin{figure}[H]\centering\includegraphics[width=0.8\linewidth]{scaling.pdf}
% \caption{Aggregate throughput vs.\ worker count.}\label{fig:scaling}\end{figure}
\todo{optional: generate the scaling figure from notebooks/graphing.ipynb and
interpret the saturation point}
